\section{개요}\label{sec:overview}

\subsection{목적}\label{subsec:purpose}

조직은 엔드포인트 탐지·대응(EDR), 네트워크 탐지·대응(NDR), 보안 오케스트레이션·자동화·대응(SOAR) 같은 방어 제품을 다층적으로 도입하지만, 이러한 제품의 실제 공격 행위에 대한 차단 및 탐지 성능은 운영 과정에서 확인하기 어렵다.
설정 문서나 제조사 자료에 명시된 탐지 가능 여부와 실제 공격 기법 실행 시의 경보 발생 여부는 서로 다를 수 있기 때문이다.

침해 및 공격 시뮬레이션(Breach and Attack Simulation, BAS)은 이러한 차이를 실제 시험을 통해 검증하는 방법이다.
알려진 공격 기법을 통제된 방식으로 대상 호스트에서 실제로 실행하고, 그 결과와 방어 제품이 남긴 기록을 비교하여 방어 체계의 취약 지점을 식별한다.

본 문서는 이러한 시뮬레이션을 폐쇄망 안에서 수행하기 위한 \textbf{BAS 플랫폼}(이하 ``플랫폼'')의 개념과 구조를 기술한다.
플랫폼은 다음 세 가지를 목표로 한다.

\begin{itemize}
    \item \textbf{실제 실행}: 공격을 모사한 트래픽이나 로그만 생성하는 것이 아니라, 대상 호스트에 상주하는 에이전트가 해당 공격 기법의 명령을 실제로 실행한다.
    \item \textbf{표준 체계 연계}: 모든 공격 절차를 MITRE ATT\&CK의 전술·기법에 대응시켜, 시험을 수행한 기법과 수행하지 않은 기법을 일관된 기준으로 구분할 수 있도록 한다.
    \item \textbf{폐쇄망 운용}: 공격 실행에 필요한 기능은 외부 서비스나 인터넷 연결 없이 동작하며, 에이전트 바이너리까지 내부에서 빌드·배포하여 시험망 내에서 독립적으로 운용한다.
\end{itemize}

\subsection{기본 개념}\label{subsec:concepts}

플랫폼의 주요 용어는 다음과 같으며, 플랫폼 화면에서 사용하는 명칭을 따른다.

\begin{itemize}
    \item \textbf{에이전트}: 대상 호스트에 상주하며 명령을 수행하는 프로그램이다. 주기적으로 C2 서버에 접속하여 명령을 수신하고, 명령을 실행한 결과를 보고한다.
    \item \textbf{C2 서버}: 에이전트와의 통신을 전담하는 서버이다. 에이전트를 등록하고, 대기 중인 명령을 전달하고, 실행 결과를 수신하여 저장한다.
    \item \textbf{Ability}: 개별 공격 절차를 정의하는 단위이다. 해당 공격 절차의 전술·기법, 지원 플랫폼, 실행기 및 실행 명령을 포함한다.
    \item \textbf{공격 실행}: 하나 이상의 Ability를 지정된 순서에 따라 대상 에이전트에서 실행하는 작업 단위이다. 단계 간 값을 전달하기 위한 \textbf{변수 저장소}를 포함한다.
    \item \textbf{평가}: 실행할 Ability와 대상 에이전트를 정의한 시험 계획이다(\ref{subsec:assessment}절).
    \item \textbf{연동}: 방어 제품의 로그 및 탐지 기록을 수집하여 시뮬레이션 결과와 비교할 수 있도록 지원하는 기능이다.
\end{itemize}

\subsection{전체 구조}\label{subsec:arch}

그림~\ref{fig:arch}은 플랫폼의 구성 요소와 각 요소 간 데이터 흐름을 나타낸다.
운영자는 웹 기반 \textbf{운영 콘솔}에서 에이전트를 배포하고, 공격 실행을 구성하며, 결과와 수집 로그를 확인한다.
대상 호스트의 에이전트는 운영 콘솔이 아닌, 별도의 프로세스로 실행되는 \textbf{C2 서버}에만 접속한다.
두 프로세스는 명령과 실행 결과 등의 상태 정보를 \textbf{공유 데이터베이스}를 통해서만 주고받는다.

\begin{figure}[H]
\centering
\begin{tikzpicture}[every node/.style={font=\footnotesize}]
  % operator side
  \node[ext, minimum width=24mm] (op) at (0,0) {운영자\\\scriptsize 웹 브라우저};
  \node[box, minimum width=30mm, minimum height=12mm] (con) at (0,-2.1) {운영 콘솔\\\scriptsize 화면 · 공격 실행 구성};
  \node[box, minimum width=30mm, minimum height=12mm] (db) at (4.4,-2.1) {공유 데이터베이스\\\scriptsize 에이전트 · 명령 · 결과};
  \node[hot, minimum width=30mm, minimum height=12mm] (c2) at (8.8,-2.1) {C2 서버\\\scriptsize 에이전트 통신 전담};
  \node[box, minimum width=30mm] (eng) at (4.4,-4.2) {공격 실행 엔진\\\scriptsize 콘솔·C2가 함께 사용};
  % target side
  \node[ext, minimum width=26mm] (a1) at (8.8,0.35) {대상 호스트 A\\\scriptsize 에이전트};
  \node[ext, minimum width=26mm] (a2) at (12.0,0.35) {대상 호스트 B\\\scriptsize 에이전트};
  \node[ext, minimum width=26mm] (sec) at (0,-4.2) {방어 제품\\\scriptsize EDR · NDR · SOAR 등};

  \draw[arr] (op) -- (con);
  \draw[arr, <->] (con) -- (db);
  \draw[arr, <->] (db) -- (c2);
  \draw[arr] (a1) -- node[lbl, right, pos=0.45]{주기적 접속} (c2);
  \draw[arr] (a2) |- (c2.east);
  \draw[arr, dashed] (eng) -- (db);
  \draw[arr, <->] (sec) -- node[lbl, right, align=left]{로그 조회 /\\로그 전송} (con);

  \begin{scope}[on background layer]
    \node[zone, fit=(op)(con)(db)(eng)(sec), label={[lbl]above:운영자 영역(내부망)}] {};
    \node[zone, fit=(c2)(a1)(a2), label={[lbl]above:시험망}] {};
  \end{scope}
\end{tikzpicture}
\caption{플랫폼 구성과 흐름}
\label{fig:arch}
\end{figure}

\subsection{설계 원칙}\label{subsec:principles}

\paragraph{\textbf{콘솔과 C2의 분리:}}
플랫폼은 C2 서버와 운영 콘솔을 별도의 프로세스로 구성하고 서로 다른 포트를 할당하였다.
시험망에서는 C2 서버에만 접근할 수 있도록 하고, 공격 실행을 제어하는 운영 콘솔은 내부망에서만 접근할 수 있도록 배치한다.
C2 서버는 에이전트 통신과 페이로드 목록 조회 기능만 제공하므로, C2 서버에 접근하더라도 운영 콘솔의 기능에는 직접 접근할 수 없다.
단, 시험망의 방어 제품에서 에이전트 방식으로 로그를 수집하는 경우에는 해당 제품 호스트에서 운영 콘솔의 수집 엔드포인트로의 접근만 추가로 허용한다(\ref{subsec:integ-token}절).

\needspace{7\baselineskip}
\paragraph{\textbf{데이터베이스를 통한 상태 공유:}}
운영 콘솔과 C2 서버는 명령과 결과를 직접 주고받지 않는다. 콘솔이 명령을 `대기' 상태로 기록하면, C2가 에이전트의 다음 정기 비컨에서 이를 전달하고, 에이전트가 보고한 결과를 다시 기록한다.
콘솔의 C2 서버 직접 호출은 공격 실행 구성 시 페이로드 파일 목록을 조회하는 용도로 한정되며(\ref{subsec:facts}절), C2 서버는 콘솔을 호출하지 않는다.
이 방식은 두 프로세스를 독립적으로 배포하거나 재시작할 수 있도록 하며, 데이터베이스에 저장된 명령과 실행 결과는 각 프로세스의 중단 및 재시작 시에도 유지된다.
공격 실행의 진행 규칙(\ref{sec:operations}장)은 콘솔과 C2가 공유하는 엔진에 구현되어 있어, 콘솔에서 공격 실행을 시작할 때와 C2에서 결과를 받을 때 같은 규칙이 적용된다.

\paragraph{\textbf{Pull 기반 명령 전달:}}
서버가 에이전트에 접속하는 것이 아니라 에이전트가 서버에 접속하여 대기 중인 명령을 수신한다.
대상 호스트가 방화벽이나 NAT 뒤에 있더라도 C2 서버로 향하는 연결이 허용되면 동작하며, 이는 실제 공격 도구에서 사용되는 통신 방식 중 하나이다.
다만 명령 전달은 에이전트의 접속 주기만큼 지연될 수 있으며, 이를 줄이는 방법은 \ref{subsec:terminal}절에서 설명한다.

\paragraph{\textbf{폐쇄망 독립 운용:}}
에이전트는 내부에 보관된 소스코드를 기반으로 운영체제(Linux, Windows, macOS)와 아키텍처(x86-64, 64비트·32비트 ARM)별로 사전 빌드되며, 생성된 바이너리는 C2 서버에 탑재된다.
외부 C2 프레임워크 없이 동작하며, 인터넷 연결이 가능한 환경에서 설치 패키지를 생성한 후 폐쇄망으로 반입하는 절차를 제공한다.
다만 외부 클라우드 서비스(예: VirusTotal)와의 API 연동은 해당 서비스에 접속할 수 있어야 하므로, 폐쇄망에서는 내부에 설치된 제품만 연동할 수 있다(\ref{subsec:integ-modes}절).

\subsection{문서 구성}

\ref{sec:c2}장은 에이전트의 배포와 C2 통신, 실시간 터미널을 다룬다.
\ref{sec:abilities}장은 공격 절차 라이브러리와 ATT\&CK 커버리지 산정 방식을, \ref{sec:operations}장은 다중 에이전트를 지원하는 공격 실행 엔진과 평가 계획을 설명한다.
\ref{sec:integrations}장은 방어 제품 연동을 설명한다.

본 문서에 수록된 화면은 모두 실제 플랫폼에서 캡처하였다.
시험 대상 호스트로는 플랫폼이 설치된 Linux 서버에서 실행한 컨테이너를 사용했으며, 해당 컨테이너에 에이전트 1개를 배포하여 무해한 정보 조회 명령만 수행하였다.
